\section{ReasonLink Design}
Figure~\ref{fig:pipeline} shows the shared-memory data paths and control interface.
\begin{figure*}[t]
\centering
\begin{tikzpicture}[font=\small,>=Latex,
 box/.style={draw,rounded corners,align=center,minimum height=11mm,text width=26mm},
 pathlabel/.style={font=\scriptsize,align=center,fill=white,inner sep=2pt}]
\node[box,fill=blue!8] (nn) at (0,0) {NVDLA tile\\int8 logits};
\node[box,fill=gray!10,text width=32mm] (dram) at (4,0) {Coherent DRAM\\logits + base + delta};
\node[box,fill=orange!12] (rl) at (8,0) {ReasonLink\\hardware-generated delta};
\node[box,fill=green!10] (sat) at (12,0) {SAT-Accel\\load + search};
\node[box,fill=gray!8,text width=44mm] (cpu) at (4,1.9) {Ariane CPU / Linux\\full image / CPU-generated delta\\descriptor and credit control};
\draw[->,thick] (nn) -- node[pathlabel,above]{output} (dram);
\draw[->,thick] (dram) -- node[pathlabel,above]{logit reads} (rl);
\draw[->,thick] (rl.south) -- ++(0,-0.65) -| node[pathlabel,pos=0.25,below]{hardware-generated delta commit} (dram.south);
\draw[->,thick] (dram.south east) -- ++(0,-1.35) -| node[pathlabel,pos=0.45,below]{solver image / delta reads} (sat.south);
\draw[<->,thick] (cpu) -- node[pathlabel,right]{CPU-generated update} (dram);
\draw[->,dashed] (cpu.east) -| node[pathlabel,pos=0.4,above]{descriptors / credits} (rl.north);
\node[font=\scriptsize,align=center,text width=30mm] at (12,1.6) {Optional on-chip base reuse:\\retain structure;\\restore mutable state};
\end{tikzpicture}

\caption{Implemented data paths. NVDLA and SAT-Accel occupy separate ESP tiles. The CPU baseline reads logits and updates a solver image in shared DRAM. ReasonLink receives CPU descriptors, reads logits coherently, queues literals on chip, and commits a compact delta to DRAM. SAT-Accel optionally retains the structural base on chip across solves. Arrows show logical data movement; all cross-tile data traverses the ESP memory system.}
\label{fig:pipeline}
\end{figure*}

\subsection{Compact Fixed Decisions and DRAM Base Reuse}
Both full-image baselines reconstruct the solver representation for every puzzle, including the unchanged structural constraints. The \emph{faithful full-image} baseline follows SAT-Accel's original host-side clearing and transfer policy: it clears capacity-sized clause and literal buffers and transfers allocated image regions, including unused capacity. The \emph{lean full-image} baseline clears previously used clause and literal regions and transfers only the currently used prefixes of the image buffers. On the board, these transfer policies move approximately 5.29\,MB and 632\,KB per puzzle, respectively. The faithful volume is half that of SAT-Accel's default build (10.5\,MB) because clause and literal-store capacities are halved to fit the VCU118 (Section~\ref{sec:impl}). Lean preparation removes unused-capacity overhead but still rebuilds and transfers the complete used representation for every puzzle. Delta updates additionally reuse the existing base image and write only the per-puzzle update.

Both delta paths initialize the base image once and update fixed decisions and counters for each puzzle. In the \emph{CPU-generated delta} path, CPU software reads logits, computes argmax, encodes fixed-decision literals, and writes the update. In the \emph{hardware-generated delta} path (ReasonLink), hardware performs these operations under CPU descriptor control. Both produce equivalent solver inputs, with different transfer layouts. \emph{On-chip base reuse} is a separate solver optimization available to either path. Across the original 1,000-puzzle dataset, only approximately 1,212 bytes of a roughly 632\,KB used image change on average. The implemented CPU-generated delta path moves 3,940 bytes per puzzle, including its update layout and bookkeeping. This is a protocol transfer volume, distinct from the number of semantically changed bytes.

\subsection{Hardware-Generated Delta (ReasonLink)}
ReasonLink resides in the SAT-Accel tile. For each given, software supplies the cell index and the address of NVDLA's output. The core reads ten signed int8 logits from coherent memory, computes argmax over digits 1--9, maps the result to a literal, and queues it in an on-chip FIFO. A commit operation writes the answer-stack delta and relevant counters. The current board path moves 2,948 bytes per puzzle in this commit protocol.

NVDLA reuses an output buffer, so the CPU waits for consumption credit before permitting reuse. This prevents the producer from overwriting logits that ReasonLink has not read. Only the core's conversion of a cell, approximately 85 tile cycles (about 1\,$\mu$s) on the board, can overlap the next inference; descriptor submission and credit polling remain CPU work and cost considerably more (Section~\ref{sec:eval}). The design thus evaluates hardware delta generation behind a shared-memory interface rather than a direct accelerator-to-accelerator wire. The project proposal envisioned a FIFO that bypasses DRAM and a solver that loads state as literals arrive. Because NVDLA and SAT-Accel occupy separate ESP tiles, literals instead reach SAT-Accel through a DRAM commit, and search begins only after the complete delta is committed.

% TODO (for later): \draftnote{Sweep FIFO depth and delta-transfer granularity with fixed predictions and a controlled producer rate. Measure backpressure, CPU descriptor/credit cost, and hardware utilization.}

\subsection{On-Chip Base Reuse and State Restoration}
DRAM base reuse removes host reconstruction but still permits SAT-Accel to reload the same base into on-chip arrays each solve. Resident-base execution additionally retains the base clause store, clause metadata, and literal occurrence storage on chip. A fresh solve populates these arrays; subsequent solves with the same base skip their copy-in.

Correct reuse requires restoring the pre-solve invariant. Base clauses are immutable, but learned clauses can occupy the free tail of each literal's final base page and link to additional pages. Before the next search, the resident path clears those tails and their links. Learned-clause allocators are reinitialized. Mutable clause states, variable metadata, and the answer stack are still copied from DRAM; this implementation does not retain and restore every state structure on chip. Sudoku copy-in falls from roughly 632\,KB to 145\,KB. The tail cleanup is a single line write per literal polarity (1,458 for Sudoku).

The driver enables residency only after a fresh execution with the same base. Back-to-back puzzles must produce the same answers and search counters as independent full-copy executions. Functional probes also compare retained base regions byte for byte after restoration. Merely checking the final Sudoku grid would miss latent learned-state contamination.

% TODO (for later): \draftnote{Add a mechanism-level control/restoration diagram, restoration pseudocode, and state invariants, including learned-clause cleanup, base changes, and buffer reuse. Relate these mechanisms to the incremental-solving techniques discussed in related work.}

\subsection{Execution Overlap}
Serial execution uses one buffer set and completes perception, handoff, and reasoning for puzzle $i$ before starting puzzle $i+1$. Double-buffered execution alternates between two buffer sets: inference, conversion, and transfer for puzzle $i+1$ proceed while SAT-Accel solves puzzle $i$. If that solve is still running when preparation finishes, execution waits for completion before starting the next solve. This scheduling choice is independent of ReasonLink's on-chip literal FIFO, which is used in both execution modes. The solver is not allowed to consume an incompletely committed delta. Base residency remains valid across the two buffers because they represent the same structural formula; retained clause identifiers and on-chip indices do not depend on a particular DRAM buffer address.
